\section{Provenance and contribution disclosure}\label{app:provenance}

This first manuscript version was assembled on 19 September 2026
by OpenAI GPT-6 (Codex), for Edward Baker, from the Wilson--Loewner
investigation and its parent research. It is a working synthesis,
not a submitted or independently refereed paper. The title-page
designation records for whom it was drafted and does not assert
human verification of every derivation. No independent review of
this package has yet been completed.

The principal local source notes, all dated 19 September 2026, are
\begin{enumerate}
 \item \emph{Smooth Wilson variation and the constant Loewner driver},
 supplying Sections~\ref{sec:geometry} and \ref{sec:variation};
 \item \emph{Defect-compatible contour freedom and the endpoint pairing},
 supplying Sections~\ref{sec:projectors} and \ref{sec:endpoints};
 \item \emph{Physical reflection, the reference junction, and a first
 bulk response}, supplying Sections~\ref{sec:reflection} and
 \ref{sec:response}.
\end{enumerate}
They are retained with the earlier handoffs in the live investigation's
\texttt{notes/} directory. The parent Wilson-lines continuation supplies
the endpoint motivation, free angular regulator and auxiliary boundary
Gaussian; the Loewner and fractional-dimension manuscripts supply
the arithmetic factorization and the scoped realization obstructions.
All formulas needed for the present arguments have been included
here. The manuscript has no shared-background TeX dependency.

The derivations use published conventions where stated. The printed
equations for the Loewner normalization, tangent-coupled projectors,
D3--D5 projector, Baker endpoint transformations and semicircle
control, Dorn--Pershin reflection rules, and the Weil-form weights
were checked against the corresponding PDFs. Bibliographic metadata
and primary links were checked during preparation. This is a focused
background review, not an exhaustive literature or priority survey.
No novelty is claimed for standard Wilson variation, tangent-coupled
supersymmetry, the ray endpoint selection rule, or the classical
Weil criterion.

The machine-readable \nolinkurl{MANUSCRIPT_SOURCES.json} identifies
the local research inputs and their hashes at assembly. The
\nolinkurl{BUILD_RECORD.json} binds the reviewed PDF to its TeX
sources and supporting checks; \nolinkurl{SNAPSHOT.json} binds a
dated archive. Neither record is a proof certificate. Historical
notes and snapshots are preserved so that a later correction can
be distinguished from the claims made in this version.
